\part{Beyond Individual Interpretation}

\chapter{Shared Models, Institutions, and Collective Self-Confirmation}

The argument so far has treated interpretation as the achievement of an individual organism, but the organization of experience is largely inherited and socially maintained. Language supplies the distinctions that a person can draw before any inquiry begins. Education trains attention and fixes which discrepancies a learner is taught to notice. Scientific instruments, record-keeping practices, habits and institutions determine what can be encountered at all and what can be recognized as evidence when it is. A model that an individual holds is therefore frequently a local instance of a model that a community maintains, and the questions of the earlier chapters, about insulation, precision and revisability, recur at the collective scale with their own mechanisms.

The formal treatment introduces coupled agents $i=1,\dots,n$ with distinct observations, models and preferences, who exchange reports and influence one another's updates. The simplest tractable case is linear averaging. Let $x_t\in\mathbb{R}^n$ collect the agents' current estimates of a quantity and let $W$ be a row-stochastic matrix recording how much weight each agent places on each other's estimate, so that
\[
x_{t+1}=Wx_t .
\]
When the network is connected and aperiodic, the estimates converge to a consensus, and the consensus value is $\pi^{\mathsf T}x_0$, where $\pi$ is the left eigenvector of $W$ for eigenvalue one, normalized to sum to one. The weights $\pi_i$ measure the long-run influence of each agent's initial estimate. Agreement is thus guaranteed by the structure of the coupling, but its content depends on who holds influence. Golub and Jackson showed that if initial estimates are independent noisy signals of a truth, the consensus converges to the truth as the group grows only if no agent retains non-vanishing influence \citep{golub2010}. Where one authority is heavily weighted, the group converges just as firmly, and converges on that authority's initial error.

Agreement can arise in several ways that the observation of agreement alone does not distinguish. It can arise from shared evidence, in which case the convergence is justified to the extent that the evidence is good. It can arise from dependence on a common authority, in which case a single source is counted many times. It can arise from coordinated inquiry in which agents test one another's claims against independent encounters. And it can arise from the suppression of dissent, in which disagreement is removed from the data of the system and not resolved. The double-counting case admits a simple statement. Suppose that $n$ agents each form a log-odds estimate $\ell$ for a hypothesis from one and the same signal, and that a naive pooling procedure sums the reports as though they were independent. The pooled value is $n\ell$, although the evidence supports only $\ell$, so that apparent confidence grows with the number of agents while the information stays fixed. A community in which agents cite one another, each citing in turn the same source, displays exactly this structure.

The notion of insulation introduced in Chapter~6 carries over. Let $\mu$ denote the distribution over situations that a community's institutional practices lead it to encounter, and $\nu$ the distribution over situations that matter for its continued viability. A large gap between the error the community's model makes on $\nu$ and on $\mu$ is compatible with collective selective exposure, a self-confirming arrangement in which instruments, categories and incentives combine to supply only the evidence the model already handles. Like its individual counterpart, the gap may instead reflect legitimate specialization, and which reading is correct requires evidence about how the encountered situations were selected. Where the arrangement is one of self-confirmation, because it is distributed, no individual need intend it, and no individual may be in a position to see it. The chapter accordingly examines the conditions under which collective convergence improves knowledge and those under which it merely stabilizes a shared error. Convergence improves knowledge when the sources feeding it are independent, when influence is not concentrated, and when disagreement can enter the record. It stabilizes error when these conditions fail, and from inside the community the two cases may look alike.

\chapter{Power, Truth, and the Availability of Correction}

A perspectival account of knowledge faces an obvious objection. If every interpretation is a perspective shaped by interests and capacities, there seems to be no ground for distinguishing stronger inquiry from effective self-insulation, since both are ways of organizing experience in the service of an organization's persistence. This chapter proposes a criterion and argues for it. The criterion concerns an interpretation's ability to sustain consequential encounters with what it does not already accommodate. Such a criterion requires argument, and it is not presented as an automatic consequence of either Nietzsche or active inference.

Nietzsche supplies a resource in the third essay of the \emph{Genealogy}, where he rejects a view from nowhere and claims that the more affects are allowed to speak about a thing, and the more different eyes can be brought to observe it, the more complete the concept of that thing, and the objectivity attainable, will be. Perspectivism in this form is a claim about the conditions of better knowing: knowledge improves with the number and difference of the perspectives brought to bear and not through escape from perspective. The remark already points toward the institutional dimension of the present chapter, since a plurality of perspectives has to be housed somewhere and its challenges have to be able to reach the organization they challenge.

The discussion connects epistemic authority to the practical availability of correction, and the following definition states the three conditions that the criterion involves.

\begin{definition}[Correctability]
An organization is \emph{correctable} with respect to a class of challenges $\mathcal{C}$ if, for each $c\in\mathcal{C}$, (i) the challenge can be registered by the organization at a precision that is not negligible, (ii) the cost borne by the party raising it is bounded and does not exceed what that party can reasonably sustain, and (iii) a revision in response to the challenge does not destroy the conditions of the organization's own activity.
\end{definition}

Each condition names a way correction fails. The first fails when precision on the relevant channel has been driven to zero by the mechanisms of Chapter~7, so that the challenge is received and not heard. The second fails when raising a challenge exposes the challenger to penalties severe enough that challenges are not raised, which makes the record of dissent a poor guide to the existence of grounds for it. The third fails when the organization cannot revise without collapse, so that even a challenge heard and welcomed cannot be acted on, and it is the condition through which the problem of transition cost from Chapter~9 enters. Who can challenge a model, what costs do they bear, and can the organization receiving the challenge change without destroying the conditions of its own activity are, in this formulation, three questions with three distinct answers.

This gives the monograph an institutional dimension and prevents self-overcoming from remaining an exclusively heroic individual achievement. The supports discussed in Chapter~11 for a person reappear as the arrangements through which a community can hear and absorb correction: procedures for adversarial testing, protected channels for dissent, redundancy of independent sources, and structures that spread the cost of revision. Where such arrangements exist, a community can transform itself without the intervention of an exceptional individual. Where they are absent, a community depends on the rare person willing to bear the full cost of challenge, and the conditions of the third kind make it likely that even then nothing changes.

The criterion should not be mistaken for a theory of truth, and it carries a limit that the final chapter records. It identifies the availability of correction as a mark of an inquiry that is capable of becoming better, and it is silent on whether the direction of change is an enhancement. An organization can be highly correctable and be corrected toward a worse condition by challenges that are themselves illusory. The criterion thus discriminates between self-insulating and self-revising organizations, which is what the perspectival account needed, without deciding which revisions deserve affirmation.

\chapter{The Limits of the Synthesis}

A comparison should state where it ceases to hold. The preceding chapters identified several points of real contact between Nietzsche's treatment of interpretation and the active-inference family of models. Schematization corresponds to coarse-graining relative to a task, with the difference that the formal notion is exact about what is lost. The ego as an interpretation of activity corresponds to a self-model that summarizes an organism's processes without disclosing them. The dependence of an error's influence on its weighting corresponds to precision, and the dependence of evidence on what the organism does corresponds to the policy-dependence of observations. Each of these correspondences is a resemblance of problem, supported by a formal construction, and each was advanced with a statement of what it leaves out.

Three limits deserve final statement. The first is whether will to power can be represented through a scalar objective. The argument of Chapter~4 showed that a weighted sum of objectives is a coherent form only where the objectives are commensurable, and Chapter~12 argued that the will to power functions in Nietzsche's use as an organizing activity within which preferences are formed, and not as a preference of high weight. If so, the attempt to identify it with the quantity minimized in active inference, or with any single information-theoretic measure, mistakes the category of the notion. The second limit concerns evaluative conflicts, which may not admit a common measure at all. Where the conflict is between ways of life and not between values on a shared scale, the formalism's tools of weighting and tradeoff presuppose what is in dispute. The third limit is whether literary transformation can be captured by changes in a generative model. Chapters~13 to 15 found that the predictive vocabulary describes some of what a text does to a reader, namely the shifting of expectation and the reallocation of attention, but that it needs added structure to represent irony and allusion, and that it offers no way of separating productive plurality from idle uncertainty or of saying why one disruption illuminates while another merely disorients.

A distinction runs through all three limits and should be kept explicit in the conclusion. Empirical questions about cognitive mechanisms, such as whether perception is implemented by hierarchical prediction, are answered by evidence about organisms, and the formal tools may prove accurate or inaccurate as models of them. Philosophical judgments about which forms of life deserve affirmation are questions of a different kind, and no result about mechanism settles them. The syntheses most likely to mislead are those that present a finding of the first kind as though it carried a verdict of the second, or that borrow the authority of a mathematical result for a valuation that the mathematics does not contain.

The conclusion that the book defends preserves a productive asymmetry. Active inference offers formal tools for understanding perception, action, learning and exploration, and it supplies exact vocabulary for what it is to preserve an organization and to revise a model within a family. Nietzsche presses questions about the authority and transformation of the organizations those tools describe, in particular how a family of models or a hierarchy of preferences acquires its standing and what it would be to replace it. The conjunction is most useful when each exposes questions left open by the other. From the formal side, the open questions are the origin of structural novelty and the authority of the objective. From Nietzsche's side, they are the conditions under which self-overcoming can be sustained and the standards by which its results are to be judged, and the chapters on attachment, institutions and correctability suggest that part of the answer lies in the relations and arrangements that make revision bearable.

The question posed at the outset was how an organism can remain sufficiently organized to act while retaining the capacity to revise the interpretations through which action becomes possible. The answer developed here is conditional and in part negative. The capacity to revise depends on exposure that is not engineered to confirm, on weights that can themselves be challenged, on supports that distribute the cost of transformation, and on access to alternatives that lie outside the family currently in use, which optimization within that family does not itself explain. Undergoing such revision may therefore require more than accepting a better proposition. It may require changes in attention, affect, identification and the reader's relation to authority, which is why a book about interpretation had to attend to how its principal author writes, and not only to what he says.
